\chapter{Lock-Free Programming}\label{ch:ll:lock-free-programming}

In July 1997 a spacecraft on the surface of Mars began resetting itself, losing data each time. A low-priority task that
gathered weather data held a lock that a high-priority task managing the information bus needed; a medium-priority task
ran in between, the high-priority task missed its deadline, and a watchdog reset the machine. The fix, uploaded from
Earth, switched on priority inheritance for that lock (Jones, 1997). A trading thread that shares a lock with a logger,
a risk reporter or a configuration reloader has the same problem, only the deadline is a few microseconds, and the lower
thread need not be preempted by anything more sinister than the operating system's scheduler. This chapter is about
sharing data between threads without locks: what the hardware guarantees about memory, what C++ and Rust let a program
say about ordering, the \omterm{def:ll:lock-free-programming:cas}{compare-and-swap} loop and its classic trap, the reclamation of memory that other threads may
still be reading, and what ``\omterm{def:ll:lock-free-programming:progress}{lock-free}'' actually promises.

\section{Atomics and the memory model}

\begin{definition}[Atomic operation]\label{def:ll:lock-free-programming:atomic}
An \emph{atomic operation}\index{atomic operation} on a memory location is indivisible with respect to other threads: no
thread observes it half done. C++ provides \texttt{std::atomic<T>} and Rust \texttt{AtomicU64} and its siblings, with
loads, stores, exchanges, fetch-and-add and \omterm{def:ll:lock-free-programming:cas}{compare-and-swap}; on x86 each maps to one instruction for sizes up to eight
bytes.
\end{definition}

A \omterm{def:ll:rust-for-low-latency:race}{data race} (chapter~9) is \omterm{def:ll:cpp-for-latency-iii-the-compiler:ub}{undefined behaviour} in C++; atomics are how a program shares a location without one. The
question they leave open is order. A processor executes out of order (chapter~2), keeps stores in a buffer before they
reach the cache, and a compiler reorders memory accesses that it can prove independent in one thread; another thread can
therefore see a thread's writes in a different order from the program's.

\begin{definition}[Memory ordering, acquire--release ordering, sequential consistency]\label{def:ll:lock-free-programming:ordering}
The \emph{memory ordering}\index{memory ordering} of an \omterm{def:ll:lock-free-programming:atomic}{atomic operation} states what other memory accesses it orders.
In \emph{acquire--release ordering}\index{acquire--release ordering}, a load with acquire that reads the value written by a
store with release sees every write the storing thread made before that store: the pair publishes data. \emph{Sequential
consistency}\index{sequential consistency}, the default of C++ atomics, adds a single total order of all such operations
that every thread agrees on, as if they were interleaved on one processor. A relaxed operation is atomic but orders
nothing else.
\end{definition}

The publication pattern (write the data, then store a flag with release; load the flag with acquire, then read the data) is
the backbone of every ring buffer in chapter~12. \omterm{def:ll:lock-free-programming:ordering}{Sequential consistency} is needed less often and costs more; the classic
case where acquire--release is not enough is the store-buffer pattern of \cref{fig:ll:lock-free-programming:litmus}: each
thread writes its own flag and reads the other's. On x86 a store may still sit in the processor's store buffer when the
following load of another location completes, so both threads can read zero. Only a sequentially consistent store (compiled
to an instruction that drains the buffer) forbids it.

\begin{omfigure}
\begin{tikzpicture}
\begin{axis}[width=9cm,height=5cm,ybar,bar width=16pt,ymin=0,enlarge y limits={upper,value=0.2},
  symbolic x coords={relaxed,release-acquire,seq_cst},xtick=data,xticklabels={relaxed,{release--acquire},seq\_cst},
  enlarge x limits=0.3,ylabel={$r_1 = r_2 = 0$ per million},tick label style={font=\small},label style={font=\small},
  grid=major,grid style={gray!25},nodes near coords,nodes near coords style={font=\footnotesize},table/col sep=comma]
\addplot[fill=omThm!45,draw=omThm] table[x=ordering,y=both_zero]{figdata/low-latency/11-lock-free-programming/measured_litmus.csv};
\end{axis}
\end{tikzpicture}

\omcaption{The store-buffer litmus test: thread A runs \texttt{x = 1; r1 = y}, thread B runs \texttt{y = 1; r2 = x},
a million times each. Under \omterm{def:ll:lock-free-programming:ordering}{sequential consistency} at least one read sees 1; with weaker orderings, the stores still in each
core's store buffer let both read 0, rarely but really. Measured on a laptop (Intel Core Ultra 7 155H) under WSL2, no
isolated cores. Data: \texttt{bench\_lockfree.py}.}\label{fig:ll:lock-free-programming:litmus}
\end{omfigure}

\section{Compare-and-swap loops and the ABA problem}

\begin{definition}[Compare-and-swap]\label{def:ll:lock-free-programming:cas}
\emph{Compare-and-swap}\index{compare-and-swap} (CAS) atomically replaces the value of a location with a new value if and
only if it still equals an expected value, and reports whether it did. A CAS loop reads the current value, computes a new
one, and retries the CAS until no other thread has changed the location in between.
\end{definition}

Every \omterm{def:ll:lock-free-programming:progress}{lock-free} structure is built from CAS loops. A stack pops by reading the head node and its successor and swinging
the head from the node to its successor with a CAS; the CAS fails, and the loop retries, if another thread moved the head
in the meantime. The failure it cannot see is the change that was undone.

\begin{definition}[ABA problem]\label{def:ll:lock-free-programming:aba}
The \emph{ABA problem}\index{ABA problem} is the failure of a \omterm{def:ll:lock-free-programming:cas}{compare-and-swap} to detect that a location changed from
$A$ to $B$ and back to $A$ between the read and the CAS: the comparison succeeds although the state it was computed from
(here, $A$'s successor) is stale.
\end{definition}

\begin{omfigure}
\begin{tikzpicture}[font=\footnotesize,
  nd/.style={draw=omDef,fill=omDef!10,circle,minimum size=0.6cm,inner sep=0pt},
  fr/.style={draw=gray,fill=gray!15,circle,minimum size=0.6cm,inner sep=0pt}]
\node[anchor=south,font=\footnotesize\bfseries,align=center] at (0.9,1.9) {1. thread 1 reads\\ head and next};
\node[nd] (a1) at (0,1.4) {A}; \node[nd] (b1) at (0.9,1.4) {B}; \node[nd] (c1) at (1.8,1.4) {C};
\draw[-Stealth] (a1) -- (b1); \draw[-Stealth] (b1) -- (c1);
\node[anchor=north,align=center] at (0.9,0.9) {head = A,\\ next = B};
\node[anchor=south,font=\footnotesize\bfseries,align=center] at (5.3,1.9) {2. thread 2: pop A,\\ pop B, push A};
\node[nd] (a2) at (4.5,1.4) {A}; \node[nd] (c2) at (5.4,1.4) {C}; \node[fr] (b2) at (6.6,1.4) {B};
\draw[-Stealth] (a2) -- (c2);
\node[anchor=north,align=center] at (5.5,0.9) {head = A again;\\ B back in the pool};
\node[anchor=south,font=\footnotesize\bfseries,align=center] at (10.0,1.9) {3. thread 1:\\ CAS(head, A $\to$ B)};
\node[fr,draw=omThm,fill=omThm!15] (b3) at (9.5,1.4) {B}; \node[nd] (c3) at (10.6,1.4) {C};
\draw[-Stealth,omThm] (b3) -- (c3);
\node[anchor=north,align=center,omThm] at (10.05,0.9) {succeeds: the head is\\ a node in the free pool};
\end{tikzpicture}

\omcaption{The \omterm{def:ll:lock-free-programming:aba}{ABA problem} on a \omterm{def:ll:lock-free-programming:progress}{lock-free} stack. The head is A both times thread 1 looks, so its \omterm{def:ll:lock-free-programming:cas}{compare-and-swap} succeeds
and installs B, which thread 2 has already removed. A tag incremented at every update makes the second A compare unequal.
Data: \texttt{ll\_aba.aba}, a deterministic replay.}\label{fig:ll:lock-free-programming:aba}
\end{omfigure}

The standard cures change what is compared. A tag (a version counter) stored next to the pointer and incremented on every
update makes the three states $A$, $B$, $A$ compare as $(A,1)$, $(B,2)$, $(A,3)$. On x86-64 a pointer and a counter fit a
16-byte \omterm{def:ll:lock-free-programming:cas}{compare-and-swap}; the tutorial's stack avoids the wide instruction by storing nodes in a fixed array and packing a
32-bit index with a 32-bit tag into one 64-bit word. Its nodes are never freed, which removes the other half of the
problem.

\section{Safe memory reclamation}

\begin{definition}[Hazard pointer, epoch-based reclamation]\label{def:ll:lock-free-programming:reclaim}
In a \omterm{def:ll:lock-free-programming:progress}{lock-free} structure whose nodes are freed, a node removed by one thread may still be read by another that loaded its
address before the removal. A \emph{hazard pointer}\index{hazard pointer} is a per-thread published pointer to a node the
thread is about to read; a node is freed only when no hazard pointer points to it (Michael, 2004). In
\emph{epoch-based reclamation}\index{epoch-based reclamation} threads announce entering and leaving a global epoch; a node
removed in epoch $e$ is freed once every thread has left $e$ (Fraser, 2004).
\end{definition}

Both defer freeing until no reader can hold the node, and both put a cost on the read path: \omterm{def:ll:lock-free-programming:reclaim}{hazard pointers} a store and a
fence per node read, epochs a store on entering and leaving each operation. On a hot path the simplest reclamation is none:
fixed pools whose slots are recycled but never returned to the system (chapter~6), bounded rings whose cells are reused by
position (chapter~12), and the tagged indices of the tutorial. Reclamation matters for unbounded structures shared between
many threads, which a hot path should not have.

\section{Progress guarantees}

\begin{definition}[Lock-free, wait-free]\label{def:ll:lock-free-programming:progress}
An algorithm is \emph{lock-free}\index{lock-free} if, whatever the scheduling of the threads, some thread completes its
operation in a finite number of its own steps: a thread that stalls cannot stop the others. It is \emph{wait-free}\index{wait-free}
if every thread completes its operation in a bounded number of its own steps, whatever the others do (Herlihy, 1991).
\end{definition}

\begin{definition}[Priority inversion]\label{def:ll:lock-free-programming:inversion}
\emph{Priority inversion}\index{priority inversion} occurs when a high-priority thread waits for a resource held by a
low-priority thread that cannot run because medium-priority threads take its processor; the high-priority thread is
effectively demoted below the medium ones.
\end{definition}

The guarantees are about stalls, not speed. A \omterm{def:ll:lock-free-programming:progress}{lock-free} algorithm's threads can still collide and retry, and under
contention a CAS loop can be slower than a lock (\cref{fig:ll:lock-free-programming:counters}); what it removes is the
possibility that a preempted thread holding a lock stops everyone, and with it \omterm{def:ll:lock-free-programming:inversion}{priority inversion}. Many useful structures
sit between the categories: the bounded multi-producer queue of the build uses one \omterm{def:ll:lock-free-programming:cas}{compare-and-swap} per operation and no
locks, yet its author notes that it is not \omterm{def:ll:lock-free-programming:progress}{lock-free} in the official sense, because a producer stopped between claiming a
cell and filling it holds up the consumers of that cell. \omterm{def:ll:lock-free-programming:progress}{Wait-free} structures are rarer and usually slower on average; a
single-producer single-consumer ring (chapter~12) is \omterm{def:ll:lock-free-programming:progress}{wait-free} for free, and is the structure trading hot paths use most.

\begin{omfigure}
\begin{tikzpicture}
\begin{axis}[width=10cm,height=5.8cm,ymode=log,xmin=0.8,xmax=4.2,ymin=2,ymax=400,xtick={1,2,3,4},
  ytick={2,5,10,20,50,100,200},yticklabels={2,5,10,20,50,100,200},
  xlabel={threads incrementing},ylabel={ns per increment, per thread (log)},tick label style={font=\small},
  label style={font=\small},grid=major,grid style={gray!25},
  legend columns=4,legend style={font=\footnotesize,at={(0.5,-0.25)},anchor=north,draw=none,fill=none,
  /tikz/every even column/.append style={column sep=6pt}},table/col sep=comma]
\addplot[omThm,very thick,mark=*] table[x=threads,y=mutex]{figdata/low-latency/11-lock-free-programming/measured_counters.csv};
\addplot[omProp,thick,mark=square*] table[x=threads,y=CAS_loop]{figdata/low-latency/11-lock-free-programming/measured_counters.csv};
\addplot[omDef,thick,mark=triangle*] table[x=threads,y=fetch_add]{figdata/low-latency/11-lock-free-programming/measured_counters.csv};
\addplot[omMeth,very thick,mark=diamond*] table[x=threads,y=per_thread]{figdata/low-latency/11-lock-free-programming/measured_counters.csv};
\legend{mutex,CAS loop,fetch\_add,per-thread counters}
\end{axis}
\end{tikzpicture}

\omcaption{One to four threads counting a million events each into a shared counter protected by a mutex, updated by a CAS
loop or by \texttt{fetch\_add}, or into padded per-thread counters summed at the end. Only the design that shares nothing
scales. Measured on a laptop (Intel Core Ultra 7 155H) under WSL2, no isolated cores, threads unpinned. Data:
\texttt{bench\_lockfree.py}.}\label{fig:ll:lock-free-programming:counters}
\end{omfigure}

The counters make the chapter's second lesson concrete. With one thread every design costs a few nanoseconds. With two or
more, every design that writes one shared line pays for the line bouncing between cores (chapter~3) at each update, a CAS
loop also for its retries, a mutex also for the lock's own line; the per-thread counters, each on its own \omterm{def:ll:memory-hierarchy-and-caches:line}{cache line},
cost what they cost alone. The fastest synchronisation is the one that is not needed: give each thread its own data and
combine at the end, or hand data over in one direction only.

\section{Tutorial: a stack, a queue and a litmus test}\label{tut:ll:lock-free-programming:tut}

\begin{tutorial}
\textbf{Goal.} Build a \omterm{def:ll:lock-free-programming:progress}{lock-free} stack that is immune to ABA, measure the cost of sharing a counter, observe memory
reordering, and build a bounded multi-producer multi-consumer queue. \textbf{End state:}
\cref{fig:ll:lock-free-programming:litmus,fig:ll:lock-free-programming:counters} and a stress-tested queue in two languages.
\begin{tutsteps}
\item \textbf{A tagged stack.} The head packs a tag with a node index; push and pop are CAS loops that increment the tag.
\omcode{code/low-latency/11-lock-free-programming/cpp/ll_lockfree.hpp}{13}{43}{A Treiber stack over a fixed pool, with a tagged head.}\label{lst:ll:lock-free-programming:stack}
\item \textbf{Replay ABA.} \texttt{ll\_aba.aba(tagged=False)} replays the interleaving of \cref{fig:ll:lock-free-programming:aba}
and returns a stack whose head is a freed node; with \texttt{tagged=True} the \omterm{def:ll:lock-free-programming:cas}{compare-and-swap} fails.
\item \textbf{The queue.} A producer claims a position with one \omterm{def:ll:lock-free-programming:cas}{compare-and-swap}, writes its cell, and publishes it with a
release store of the cell's sequence number; the consumer of that position acquires it.
\omcode{code/firm/mpmcq/cpp/firm_mpmcq.hpp}{30}{48}{The producer side of the bounded MPMC queue.}\label{lst:ll:lock-free-programming:push}
\item \textbf{Stress and sanitise.} The tests run two producers and two consumers through 300\,000 items and check that each
arrives exactly once, in C++ and in Rust; the C++ test also runs under the thread sanitiser.
\item \textbf{Measure} with \texttt{python bench\_lockfree.py}.
\end{tutsteps}
\textbf{What to change next.} Remove the tag from the stack's head and run the two-thread churn test until it fails; change the
queue's release store to relaxed and run the stress test under the thread sanitiser.
\end{tutorial}

\section{Build: the bounded multi-producer multi-consumer queue}\label{bld:ll:lock-free-programming:mpmcq}

\begin{build}
\bfield{Purpose} The queue for the few places where several threads produce into, or consume from, one stream: order
requests from several strategies to one gateway, work for a pool of risk calculators. Single-producer paths use chapter~12's
rings instead.
\bfield{Interface} C++20 \texttt{firm::mpmcq::Queue<T>(capacity)} with \texttt{try\_push(v)} and \texttt{try\_pop()} returning
\texttt{std::optional<T>}; Rust \texttt{firm\_mpmcq::Queue<T: Copy>} with \texttt{try\_push} and \texttt{try\_pop}.
\bfield{Rules} Capacity a power of two; one \omterm{def:ll:lock-free-programming:cas}{compare-and-swap} per operation; no allocation after construction; full and empty
reported, never waited for inside the queue; producer and consumer positions on separate \omterm{def:ll:memory-hierarchy-and-caches:line}{cache lines}.
\bfield{Acceptance tests} \texttt{code/firm/mpmcq/}: first in, first out through a full and empty cycle; two producers and two
consumers deliver 300\,000 items exactly once, in C++20 and Rust, within a few seconds.
\bfield{Stretch} Batch operations that claim several positions with one \omterm{def:ll:lock-free-programming:cas}{compare-and-swap}; a blocking wrapper for threads that
may sleep.
\end{build}

\begin{omsources}
\item M.~Jones, ``What really happened on Mars Rover Pathfinder'', 1997.
\item H.-J.~Boehm and S.~Adve, ``Foundations of the C++ concurrency memory model'', PLDI 2008.
\item M.~Herlihy, ``Wait-free synchronization'', \emph{ACM TOPLAS} 13(1), 1991.
\item M.~Michael, ``Hazard pointers: safe memory reclamation for lock-free objects'', \emph{IEEE TPDS} 15(6), 2004;
K.~Fraser, \emph{Practical lock-freedom}, Cambridge technical report 579, 2004.
\item D.~Vyukov, ``Bounded MPMC queue'', 1024cores.
\end{omsources}

\section{Exercises}

\begin{exercise}[$\star$]\label{exo:ll:lock-free-programming:1}
From \cref{fig:ll:lock-free-programming:counters}, how much does a shared \texttt{fetch\_add} cost per increment with one
thread and with four? How many increments a second do four threads achieve together, shared and per-thread?
\end{exercise}

\begin{exercise}[$\star$]\label{exo:ll:lock-free-programming:2}
A writer stores a price, then sets a flag with release; a reader loads the flag with acquire and, if set, reads the price.
Can the reader see the flag set and an old price? And if the flag store is relaxed?
\end{exercise}

\begin{exercise}[$\star$]\label{exo:ll:lock-free-programming:3}
A tag of 32 bits is incremented at every update of the stack's head at 100 million updates a second. How long before it wraps,
and why is wrapping (almost) harmless?
\end{exercise}

\begin{exercise}[$\star\star$]\label{exo:ll:lock-free-programming:4}
In the store-buffer test, which outcomes $(r_1, r_2)$ are possible under \omterm{def:ll:lock-free-programming:ordering}{sequential consistency}, and which extra one does x86
allow with release--acquire? Why does the figure show it so rarely?
\end{exercise}

\begin{exercise}[$\star\star$]\label{exo:ll:lock-free-programming:5}
Why is the bounded MPMC queue not \omterm{def:ll:lock-free-programming:progress}{lock-free} in the strict sense? Describe the stall.
\end{exercise}

\begin{exercise}[$\star\star$]\label{exo:ll:lock-free-programming:6}
A logger thread holds a mutex for \qty{2}{\micro\second} per message, 100\,000 times a second, and the trading thread takes the
same mutex once per order. What fraction of orders find it held, and how long can one wait if the logger is preempted for a
scheduler quantum of \qty{4}{\milli\second}?
\end{exercise}

\begin{exercise}[$\star\star\star$]\label{exo:ll:lock-free-programming:7}
\emph{Coding.} Run the untagged stack of \texttt{ll\_lockfree.hpp} (mode \texttt{aba} of the benchmark) with two threads
popping and pushing back a million times each, twenty times. Is a node ever lost or duplicated? What does the result say about
testing \omterm{def:ll:lock-free-programming:progress}{lock-free} code by running it?
\end{exercise}

\begin{exercise}[$\star\star\star$]\label{exo:ll:lock-free-programming:8}
\emph{Find the flaw.} ``Our order-book snapshot is published with a relaxed flag; we tested it for a week on x86 and never saw a
torn snapshot, so it is correct.''
\end{exercise}

\section{Problem: The Descheduled Lock Holder}

\begin{problem}[{Weekend problem --- a lock shared with the logger}]\label{pb:ll:lock-free-programming:1}
A strategy thread sends about 20\,000 orders a second and, for each, takes a mutex to append the order to a log buffer that a
logger thread drains. The logger holds the same mutex for \qty{3}{\micro\second} each time it drains, 50\,000 times a second.
Once in a while the operating system preempts the logger while it holds the mutex, for a quantum of about
\qty{3}{\milli\second}; suppose this happens on 0.01\% of its drains.

\medskip\noindent\textbf{Part I --- Contention.}
\begin{enumerate}
\item What fraction of the time does the logger hold the mutex?
\item What fraction of orders find it held, and how long do they wait on average?
\item What is the expected blocking time per second for the strategy, ignoring preemption?
\item What does \cref{fig:ll:lock-free-programming:counters} say about the mutex's own cost when two threads take it often?
\end{enumerate}

\medskip\noindent\textbf{Part II --- Preemption.}
\begin{enumerate}[resume]
\item How many preemptions of the lock holder happen per second?
\item What fraction of time is the mutex held by a preempted logger?
\item How many orders a second meet such a preempted holder, and how long does each wait at most?
\item Add Parts I and II: the expected blocking per second.
\end{enumerate}

\medskip\noindent\textbf{Part III --- The redesign.}
\begin{enumerate}[resume]
\item Replace the mutex by a single-producer single-consumer ring: what does the strategy do per order now?
\item What happens when the ring is full?
\item Which progress guarantee does the ring give the strategy?
\item Why does priority inheritance not solve the problem for a trading thread?
\end{enumerate}

\medskip\noindent\textbf{Part IV --- The verdict.}
\begin{enumerate}[resume]
\item State the \emph{named result}: the strategy's expected blocking time per second with the mutex, and with the ring.
\item What did the Mars Pathfinder engineers do, and why was it enough there?
\item Where else in a trading process can a lock hide?
\item What does the \omterm{def:ll:lock-free-programming:aba}{ABA problem} require of a \omterm{def:ll:lock-free-programming:progress}{lock-free} structure that frees memory?
\item How do the build's tests check the queue, and what can they not prove?
\item When is a mutex the right choice in a trading system?
\item What \omterm{def:ll:lock-free-programming:ordering}{memory ordering} does the ring's publication need?
\item In one sentence: what does \omterm{def:ll:lock-free-programming:progress}{lock-free} buy a hot path?
\end{enumerate}
\end{problem}

\section{Interview questions}

\begin{interviewq}[$\star$ \iqroles{developer}]\label{iq:ll:lock-free-programming:1}
What is the difference between \texttt{memory\_order\_relaxed}, acquire--release and sequentially consistent atomics?
\end{interviewq}

\begin{interviewq}[$\star\star$ \iqroles{developer}]\label{iq:ll:lock-free-programming:2}
Explain the \omterm{def:ll:lock-free-programming:aba}{ABA problem} and two ways to prevent it.
\end{interviewq}

\begin{interviewq}[$\star\star$ \iqroles{developer}]\label{iq:ll:lock-free-programming:3}
Is a \omterm{def:ll:lock-free-programming:progress}{lock-free} algorithm always faster than one with locks? Give a counterexample.
\end{interviewq}

\begin{interviewq}[$\star\star$ \iqroles{developer}]\label{iq:ll:lock-free-programming:4}
What is \omterm{def:ll:lock-free-programming:inversion}{priority inversion}, and why does it matter in a trading system?
\end{interviewq}

\begin{interviewq}[$\star\star$ \iqroles{developer}]\label{iq:ll:lock-free-programming:5}
How do you test \omterm{def:ll:lock-free-programming:progress}{lock-free} code? What can a stress test not tell you?
\end{interviewq}

\begin{interviewq}[$\star\star\star$ \iqroles{developer}]\label{iq:ll:lock-free-programming:6}
Design a counter that many threads increment and one thread reads once a second, with the lowest cost to the incrementing
threads.
\end{interviewq}
